\section{Mathematical Details of the Joint Interval-DMR Model}
\label{sec:joint_interval_math}

\subsection{Observed Data and Indexing}

Let \(i=1,\ldots,N\) index patients. Patient \(i\) has \(n_i\) longitudinal molecular observations and \(K_i\) at-risk intervals for first deep molecular response (DMR). The cleaned GLW data contain \(N=87\) patients, \(\sum_i n_i=495\) longitudinal observations, and \(\sum_i K_i=275\) at-risk intervals.

For longitudinal visit \(j\) of patient \(i\), define
\[
y_{ij}=\mathrm{log\mbox{-}MRD}_{ij}, \qquad
t_{ij}=\mathrm{years\ from\ imatinib\ start}, \qquad
x_{ij}=\log(1+t_{ij}),
\]
and let \(z_{ij}=1\) if the sample type is bone marrow and \(z_{ij}=0\) otherwise. The transformed time \(x_{ij}\) is used to allow rapid early molecular change and slower later change.

For at-risk interval \(k\) of patient \(i\), define
\[
L_{ik}=\mathrm{interval\ start}, \qquad
R_{ik}=\mathrm{interval\ end}, \qquad
\Delta_{ik}=R_{ik}-L_{ik},
\]
where time is measured in years. Let
\[
t^{mid}_{ik}=\frac{L_{ik}+R_{ik}}{2}, \qquad
x^{mid}_{ik}=\log(1+t^{mid}_{ik}).
\]
The event indicator is
\[
d_{ik} =
\begin{cases}
1, & \text{if first observed DMR occurs in interval }(L_{ik},R_{ik}],\\
0, & \text{otherwise.}
\end{cases}
\]
By construction, each patient has at most one interval with \(d_{ik}=1\). Intervals before first DMR or through censoring have \(d_{ik}=0\).

\subsection{Longitudinal Molecular Submodel}

The latent patient-specific molecular trajectory is
\[
\mu_{ij}
=\beta_0+\beta_1x_{ij}+\beta_2x_{ij}^2+\beta_3z_{ij}
+b_{0i}+b_{1i}x_{ij}.
\]
Here \(\beta_0\) is the population intercept, \(\beta_1\) and \(\beta_2\) describe the population time trend, and \(\beta_3\) adjusts for sample type. The patient-specific random intercept and random slope are
\[
\mathbf{b}_i =
\begin{pmatrix}
b_{0i}\\
b_{1i}
\end{pmatrix}
\sim N_2(\mathbf{0},\Sigma_b),
\]
with
\[
\Sigma_b =
\begin{pmatrix}
\tau_0 & 0\\
0 & \tau_1
\end{pmatrix}
\Omega_b
\begin{pmatrix}
\tau_0 & 0\\
0 & \tau_1
\end{pmatrix},
\]
where \(\tau_0,\tau_1>0\) are random-effect standard deviations and \(\Omega_b\) is a \(2\times2\) correlation matrix.

For observations not reported at the assay floor,
\[
y_{ij}\mid \mu_{ij},\sigma_y \sim N(\mu_{ij},\sigma_y^2).
\]
The Stan model treats values at the reporting floor as left-censored. Let \(c=-5\) be the log-MRD floor and let \(a_{ij}=1\) indicate that the observation is at the floor. Then the longitudinal contribution is
\[
\ell^{Y}_{ij} =
\begin{cases}
\log \phi(y_{ij};\mu_{ij},\sigma_y^2), & a_{ij}=0,\\
\log \Phi\left(\dfrac{c-\mu_{ij}}{\sigma_y}\right), & a_{ij}=1,
\end{cases}
\]
where \(\phi(\cdot;\mu,\sigma^2)\) is the normal density and \(\Phi(\cdot)\) is the standard normal cumulative distribution function.

\subsection{Interval-DMR Submodel}

The event submodel links DMR onset to the latent molecular trajectory. At the midpoint of interval \(k\), the latent molecular value used by the hazard model is
\[
\mu_i(t^{mid}_{ik})
=\beta_0+\beta_1x^{mid}_{ik}+\beta_2(x^{mid}_{ik})^2
+b_{0i}+b_{1i}x^{mid}_{ik}.
\]
The current Stan implementation does not include sample type in the interval midpoint trajectory because intervals are defined at the patient-follow-up level rather than at a specific sample type.

The interval-specific log hazard is
\[
\eta_{ik}
=\gamma_0+\gamma_1x^{mid}_{ik}
+\gamma_2\log(1+\Delta_{ik})
+\alpha\,\mu_i(t^{mid}_{ik}),
\]
and the interval hazard is
\[
h_{ik}=\exp(\eta_{ik}).
\]
Assuming a piecewise-constant hazard within each observed interval, the cumulative hazard over the interval is
\[
H_{ik}=h_{ik}\Delta_{ik}.
\]
Therefore, the probability that first DMR occurs during the interval is
\[
p_{ik}
=P(d_{ik}=1\mid\mathrm{at\ risk})
=1-\exp(-H_{ik}).
\]
The interval likelihood contribution is
\[
\ell^{D}_{ik}
=d_{ik}\log\{1-\exp(-H_{ik})\}-(1-d_{ik})H_{ik}.
\]
This is exactly the likelihood form implemented in Stan:
\[
\ell^{D}_{ik} =
\begin{cases}
\log\{1-\exp(-H_{ik})\}, & d_{ik}=1,\\
-H_{ik}, & d_{ik}=0.
\end{cases}
\]

\subsection{Joint Likelihood}

Let \(\theta\) denote all fixed-effect, variance, association, and random-effect parameters. Conditional on the random effects, the joint likelihood is
\[
L(\theta)
=
\prod_{i=1}^{N}
\left[
\prod_{j=1}^{n_i} L^{Y}_{ij}(\theta)
\prod_{k=1}^{K_i} L^{D}_{ik}(\theta)
\right],
\]
or equivalently the log-likelihood is
\[
\ell(\theta)
=
\sum_{i=1}^{N}
\left[
\sum_{j=1}^{n_i}\ell^{Y}_{ij}(\theta)
+\sum_{k=1}^{K_i}\ell^{D}_{ik}(\theta)
\right].
\]
The shared random effects \(b_{0i}\) and \(b_{1i}\) induce dependence between the repeated molecular measurements and first-DMR process.

\subsection{Prior Distributions}

The Stan implementation uses weakly informative priors centered on clinically plausible molecular-response behavior:
\[
\beta_0\sim N(-2.5,2^2), \qquad
\beta_1\sim N(-1,1^2), \qquad
\beta_2\sim N(0,0.5^2),
\]
\[
\beta_3\sim N(0,1^2), \qquad
\sigma_y\sim \mathrm{Exponential}(1).
\]
For the interval-DMR submodel,
\[
\gamma_0\sim N(-2,2^2), \qquad
\gamma_1\sim N(0,1^2), \qquad
\gamma_2\sim N(0,1^2),
\]
\[
\alpha\sim N(-0.5,0.75^2).
\]
The prior mean for \(\alpha\) is negative because lower log-MRD corresponds to deeper molecular response; therefore lower molecular burden is expected to increase the probability of DMR within an interval.

For the random effects,
\[
\tau_0,\tau_1\sim \mathrm{Exponential}(1), \qquad
\Omega_b\sim \mathrm{LKJ}(2).
\]
The Stan code uses a non-centered Cholesky parameterization:
\[
\mathbf{b}_i = \mathrm{diag}(\boldsymbol{\tau})L_{\Omega}\mathbf{z}_i,
\qquad
\mathbf{z}_i\sim N_2(\mathbf{0},I_2),
\]
where \(L_{\Omega}L_{\Omega}^{T}=\Omega_b\).

\subsection{Posterior Distribution}

The posterior distribution is proportional to the product of the joint likelihood and all prior densities:
\[
p(\theta\mid Y,D)
\propto
\left[
\prod_{i=1}^{N}
\prod_{j=1}^{n_i} L^{Y}_{ij}(\theta)
\prod_{k=1}^{K_i} L^{D}_{ik}(\theta)
\right]
p(\boldsymbol{\beta})p(\boldsymbol{\gamma})p(\alpha)
p(\sigma_y)p(\boldsymbol{\tau})p(\Omega_b)
\prod_{i=1}^{N}p(\mathbf{b}_i\mid\boldsymbol{\tau},\Omega_b).
\]
Posterior inference focuses on the molecular time trend \((\beta_1,\beta_2)\), between-patient heterogeneity \((\tau_0,\tau_1,\Omega_b)\), visit-gap effect \(\gamma_2\), and association parameter \(\alpha\).

\subsection{Interpretation of Key Parameters}

\begin{itemize}
  \item \(\beta_1,\beta_2\): population-level change in log-MRD over time.
  \item \(\beta_3\): systematic difference between bone marrow and non-bone-marrow samples.
  \item \(\tau_0,\tau_1\): between-patient heterogeneity in baseline molecular level and response slope.
  \item \(\gamma_1\): change in interval-DMR hazard across follow-up time.
  \item \(\gamma_2\): association between visit-gap length and interval event probability after accounting for the interval length.
  \item \(\alpha\): association between latent molecular burden and DMR onset. A negative value supports the clinical expectation that lower log-MRD increases the chance of DMR.
\end{itemize}

\subsection{Derived Quantities}

For each interval, the generated quantity
\[
\widehat{p}_{ik}=1-\exp(-H_{ik})
\]
summarizes the posterior interval probability of first DMR. Patient-specific fitted trajectories are obtained from
\[
\widehat{\mu}_{i}(t)
=\beta_0+\beta_1\log(1+t)+\beta_2\{\log(1+t)\}^2
+b_{0i}+b_{1i}\log(1+t).
\]
These fitted trajectories can be plotted against observed log-MRD values and used to estimate individualized DMR probabilities over future monitoring intervals.

\subsection{Benchmark Models}

Because a Stan interface is not currently installed in the workspace, the pipeline also fits benchmark models using installed R packages. The interval-censored Weibull model uses
\[
T_i \sim \mathrm{Weibull}(\lambda_i,\kappa),
\qquad
\log \lambda_i = \delta_0+\delta_1 y_{i0}
\]
for the molecular-only model, with a complete-case extension
\[
\log \lambda_i
=\delta_0+\delta_1y_{i0}+\delta_2\mathrm{age}_i
+\delta_3\mathrm{male}_i+\delta_4\mathrm{duration}_i.
\]
The longitudinal benchmark model is a mixed-effects model,
\[
y_{ij}
=f(t_{ij})+\beta_3z_{ij}+u_{0i}+u_{1i}t_{ij}+\epsilon_{ij},
\]
where \(f(t_{ij})\) is a natural spline in years from imatinib start. These benchmark models verify that the cleaned data support interval-censored survival modeling and longitudinal mixed modeling, but the proposed joint model remains the primary model.

\subsection{Assumptions and Sensitivity Analyses}

The primary assumptions are:
\begin{enumerate}
  \item the latent log-MRD trajectory is adequately represented by a quadratic function of \(\log(1+t)\) plus patient-specific random effects;
  \item reported floor values at log-MRD \(\leq -5\) are left-censored observations;
  \item the DMR hazard is approximately constant within each observed monitoring interval;
  \item censoring is conditionally independent given the observed follow-up and modeled molecular trajectory;
  \item each patient can experience first DMR only once.
\end{enumerate}
Recommended sensitivity analyses are:
\begin{enumerate}
  \item fit an exact-value model that treats floor values as observed \(-5\) to assess floor-censoring impact;
  \item fit a model without \(\gamma_2\) to evaluate whether explicit visit-gap adjustment changes inference;
  \item replace the quadratic \(\log(1+t)\) trajectory with spline time terms;
  \item compare molecular-only and complete-case covariate-adjusted survival models;
  \item evaluate prior sensitivity for \(\alpha\), especially using \(N(0,1)\) and \(N(-0.5,1.5^2)\) alternatives.
\end{enumerate}

